\section{Module \texttt{date.c}}

Le module \texttt{date.c} encapsule le widget calendrier de GTK4. Il prend en charge la date initiale, l'affichage des noms de jours, les numéros de semaine et la remontée du signal de sélection.

\subsection{Fonction \texttt{create\_calendar}}

\begin{lstlisting}[language=C]
GtkWidget *create_calendar(const calendar_config *config) {
    GtkWidget *calendar = gtk_calendar_new();

    if (config == NULL) {
        return calendar;
    }

    if (config->selected_date != NULL) {
#if GTK_CHECK_VERSION(4, 20, 0)
        gtk_calendar_set_date(GTK_CALENDAR(calendar), config->selected_date);
#else
        gtk_calendar_select_day(GTK_CALENDAR(calendar), config->selected_date);
#endif
    }
    gtk_calendar_set_show_day_names(GTK_CALENDAR(calendar), config->show_day_names);
    gtk_calendar_set_show_week_numbers(GTK_CALENDAR(calendar), config->show_week_numbers);
    gtk_calendar_set_show_heading(GTK_CALENDAR(calendar), !config->no_month_change);
    apply_widget_style(calendar, &config->style);

    if (config->on_day_selected != NULL) {
        g_signal_connect(calendar, "day-selected", G_CALLBACK(config->on_day_selected), config->user_data);
    }

    return calendar;
}
\end{lstlisting}

\subsection{APIs GTK utilisées}

\begin{itemize}
    \item \textbf{\texttt{gtk\_calendar\_new()}} : Instancie le composant de base du calendrier GTK4. Cette API initialise le widget permettant l'affichage et la sélection interactive de dates dans le module.
    \item \textbf{\texttt{gtk\_calendar\_set\_date()}} : Positionne le calendrier sur une date précise (jour, mois, année) à l'aide d'un objet \texttt{GDateTime}. Disponible depuis GTK 4.20, elle est privilégiée pour définir l'état initial du widget selon la configuration.
    \item \textbf{\texttt{gtk\_calendar\_select\_day()}} : Remplit le même rôle que l'API précédente mais pour les versions de GTK4 plus anciennes. Son utilisation dans le module garantit la portabilité du code sur différents environnements système.
    \item \textbf{\texttt{gtk\_calendar\_set\_show\_week\_numbers()}} : Bascule l'affichage des colonnes de numérotation des semaines sur le côté gauche du widget. Elle est directement pilotée par le booléen \texttt{show\_week\_numbers} de la structure de configuration.
    \item \textbf{\texttt{gtk\_calendar\_set\_show\_heading()}} : Contrôle la visibilité de la barre supérieure affichant le mois et l'année. Le wrapper s'en sert pour activer ou désactiver la navigation mensuelle directe par l'utilisateur.
    \item \textbf{\texttt{g\_signal\_connect()}} : Attache un callback au signal \texttt{day-selected}. Cette connexion permet au wrapper de notifier le reste de l'application dès qu'une nouvelle date est cliquée sur le calendrier.
\end{itemize}

\newpage
